\DocumentMetadata{
  pdfversion=2.0,
  pdfstandard=ua-2,
  lang=en-US,
  tagging=on,
  tagging-setup={math/setup=mathml-SE,tabsorder=structure}
}
\documentclass[11pt,letterpaper]{article}
\usepackage[margin=0.68in,includefoot]{geometry}
\usepackage{newcomputermodern}
\usepackage{amsmath,amscd,mathtools}
\usepackage{tikz,adjustbox}
\providecommand{\square}{\mdlgwhtsquare}
\usepackage[hidelinks]{hyperref}
\hypersetup{
  pdftitle={Numerical Analysis PhD Exam, September 2000},
  pdfauthor={University of Florida Department of Mathematics},
  pdfsubject={Graduate mathematics examination},
  pdfdisplaydoctitle=true,
  bookmarksopen=true
}
\setcounter{secnumdepth}{0}
\setlength{\parindent}{0pt}
\setlength{\parskip}{0.28em}
\setlength{\emergencystretch}{3em}
\setlength{\leftmargini}{2.15em}
\setlength{\leftmarginii}{2.65em}
\setlength{\leftmarginiii}{3em}
\pagestyle{empty}
\raggedbottom
\allowdisplaybreaks
\NewTaggingSocketPlug{block/list/label}{examproblem}{%
  \tagstructbegin{tag=Lbl}\tagstructend%
  \tagstructbegin{tag=\LBody}%
  \tagstructbegin{tag=\ProblemHeadingTag,title-o={\ProblemHeadingText}}%
  \tagmcbegin{tag=\ProblemHeadingTag,actualtext={\ProblemHeadingText}}%
  #2%
  \tagmcend%
  \tagstructend%
}
\newenvironment{examproblems}[2]{%
  \def\ProblemHeadingTag{#2}%
  \AssignTaggingSocketPlug{block/list/label}{examproblem}%
  \begin{enumerate}%
  \setcounter{enumi}{#1}%
  \renewcommand{\labelenumi}{\textbf{\theenumi.}}%
  \setlength{\topsep}{0.35em}%
  \setlength{\partopsep}{0pt}%
  \setlength{\itemsep}{0.48em}%
  \setlength{\parsep}{0.18em}%
}{\end{enumerate}}
\newenvironment{examsubparts}{%
  \AssignTaggingSocketPlug{block/list/label}{default}%
  \begin{enumerate}%
  \setlength{\topsep}{0.18em}%
  \setlength{\partopsep}{0pt}%
  \setlength{\itemsep}{0.16em}%
  \setlength{\parsep}{0.1em}%
}{\end{enumerate}}
\newenvironment{examinstructions}{%
  \par\begingroup\itshape
}{%
  \par\endgroup\vspace{0.18em}
}
\makeatletter
\renewcommand\subsection{\@startsection{subsection}{2}{0pt}%
  {-1.25ex plus -.25ex minus -.1ex}%
  {0.55ex plus .1ex}%
  {\normalfont\large\bfseries}}
\renewcommand{\@maketitle}{%
  \newpage\null\vskip -1em%
  {\footnotesize\bfseries
    \href{https://www.ufl.edu/}{University of Florida}, \href{https://math.ufl.edu/}{Department of Mathematics}\par}%
  \vskip -0.5em%
  \tagstructbegin{tag=H1,title={Numerical Analysis PhD Exam, September 2000}}%
  \tagpdfparaOff%
  \tagmcbegin{tag=H1}%
    {\LARGE\bfseries\href{https://gma.math.ufl.edu/exams/numerical-analysis-phd/}{Numerical Analysis PhD Exam}, September 2000\par}%
  \tagmcend%
  \tagpdfparaOn%
  \tagstructend%
  \vskip 0.25em%
  \tagmcbegin{artifact}%
    \hrule height 1pt%
  \tagmcend%
  \vskip 1em%
}
\makeatother
\title{Numerical Analysis PhD Exam, September 2000}
\author{}
\date{}
\begin{document}
\maketitle
\thispagestyle{empty}
\begin{examinstructions}
Do any 8 of the following 10 problems.
\end{examinstructions}
\subsection{Part 1: Numerical Linear Algebra}
\begin{examproblems}{0}{H3}
\def\ProblemHeadingText{Problem 1.}
\item
A square matrix~\(A\) which can be decomposed into a form \(A=SDS^{-1}\), where~\(D\) is diagonal, is said to be diagonalizable.
\begin{examsubparts}
\item[\textnormal{(a)}]
Is the matrix
\[
A=\begin{pmatrix}5&1\\0&3\end{pmatrix}
\]
 diagonalizable? Why?
\item[\textnormal{(b)}]
Can the above matrix~\(A\) be diagonalized with a matrix~\(S\) such that \(S^*S=I\)?
\end{examsubparts}
\def\ProblemHeadingText{Problem 2.}
\item
This problem has three parts.
\begin{examsubparts}
\item[\textnormal{(a)}]
State a procedure for factoring a matrix~\(A\) (without pivoting) into a form \(A=LDU\), where~\(L\) is unit lower triangular, \(U\) is unit upper triangular, and~\(D\) is diagonal.
\item[\textnormal{(b)}]
State sufficient conditions on~\(A\) for the procedure in part (a) (without pivoting) to work.
\item[\textnormal{(c)}]
What is the computational advantage of having lower and upper triangular matrices? An explanation without specific computation counts is sufficient.
\end{examsubparts}
\def\ProblemHeadingText{Problem 3.}
\item
This problem has two parts.
\begin{examsubparts}
\item[\textnormal{(a)}]
Under what conditions can a matrix be factored into a form \(A=U^*TU\), where \(U^*U=I\), and~\(T\) is upper triangular?
\item[\textnormal{(b)}]
If~\(A\) is Hermitian, how is~\(U\) constructed? (You don’t need to actually give an algorithm.)
\end{examsubparts}
\def\ProblemHeadingText{Problem 4.}
\item
This problem has three parts.
\begin{examsubparts}
\item[\textnormal{(a)}]
State the power method algorithm for finding a dominant eigenvalue and an associated eigenvector of a matrix~\(A\). (An eigenvalue \(\lambda_i\) is said to be dominant if \(|\lambda_i|>|\lambda_k|\) for all other~\(k\).)
\item[\textnormal{(b)}]
Show that the power method converges (for almost any starting value) when~\(A\) is real symmetric and has a dominant eigenvalue.
\item[\textnormal{(c)}]
Explain how the inverse power method can be used to find non-dominant eigenvalues and eigenvectors for a matrix~\(A\).
\end{examsubparts}
\def\ProblemHeadingText{Problem 5.}
\item
This problem has three parts.
\begin{examsubparts}
\item[\textnormal{(a)}]
Give a careful statement of the Gershgorin Circle Theorem.
\item[\textnormal{(b)}]
Give a careful proof of the Gershgorin Circle Theorem.
\item[\textnormal{(c)}]
Let
\[
A=\begin{pmatrix}5&.01&.03\\.02&4&.01\\.03&-.01&4\end{pmatrix}.
\]
 What can you say about the location of the eigenvalues of~\(A\)?
\end{examsubparts}
\end{examproblems}
\subsection{Part II: Numerical Analysis}
\begin{examproblems}{5}{H3}
\def\ProblemHeadingText{Problem 6.}
\item
This problem has three parts.
\begin{examsubparts}
\item[\textnormal{(a)}]
State Newton’s method (the algorithm) for solving \(f(x)=0\), where \(f:\mathbb{R}\to\mathbb{R}\).
\item[\textnormal{(b)}]
Assuming~\(f\) is smooth and we have an initial guess sufficiently close to the root, state sufficient condition(s) for the method to converge quadratically.
\item[\textnormal{(c)}]
How would you use Newton’s method to compute \(5^{1/3}\)?
\end{examsubparts}
\def\ProblemHeadingText{Problem 7.}
\item
This problem has three parts.
\begin{examsubparts}
\item[\textnormal{(a)}]
Give an exact expression for a Lagrange polynomial of degree~\(n\) which would interpolate \(n+1\) given data points, \(f(x_k)\).
\item[\textnormal{(b)}]
Describe the differences between Hermite polynomials and Lagrange polynomials.
\item[\textnormal{(c)}]
Suppose we have a set of Lagrange polynomials of degree \(2n+1\) which interpolate a smooth function at the points \(x_0,x_0+h,x_1,x_1+h,x_2,x_2+h,\ldots,x_n,x_n+h\). What is the pointwise limit of this interpolant as \(h\to0\)?
\end{examsubparts}
\def\ProblemHeadingText{Problem 8.}
\item
This problem has two parts.
\begin{examsubparts}
\item[\textnormal{(a)}]
Describe Romberg’s technique for numerical approximation of an integral (based on the trapezoid rule).
\item[\textnormal{(b)}]
Give a brief justification of the method.
\end{examsubparts}
\def\ProblemHeadingText{Problem 9.}
\item
Triple Recursion Formula. If \(\{\phi_n(x)\}\) is an orthogonal family of polynomials on \([a,b]\), with weight function \(w(x)\geq0\), and \(n\geq1\), then show that
\[
\phi_{n+1}(x)=(a_nx+b_n)\phi_n(x)-c_n\phi_{n-1}(x),
\]
 for some constants \(a_n,b_n\), and \(c_n\). (Hint: Let \(g(x)=\phi_{n+1}(x)-a_nx\phi_n(x)\), where \(a_n\) is a constant chosen so that \(g(x)\) has degree~\(n\).)
\def\ProblemHeadingText{Problem 10.}
\item
Let \(\{\phi_k\}\) be a system of orthogonal polynomials, with respect to a non-negative weight \(w(x)\), on \([a,b]\). Show that each of the polynomials \(\phi_k(x)\) must have exactly~\(k\) simple zeros in \([a,b]\).
\end{examproblems}
\end{document}
